\documentclass[a4paper,12pt]{article}
\usepackage{amsmath}
\usepackage{graphicx}
\usepackage[margin=1in]{geometry}


\begin{document}

\title {Bijectivity and the double angle formula}
\maketitle
		

\centerline{\textbf{Injection, surjection, and bijection}}

\bigskip
Let us repeat the definitions from class. Suppose $f:X \mapsto Y$ is a map from $X$ to $Y$.
\\

\textbf{Def}: $f$ is \textbf{injective} if
\begin{center}$\forall \  a$, $b \in X$, $a \ne b \implies f(a) \ne f(b)$\end{center}
This is also called a \textbf{one-to-one} map. 

\textbf{Def}: $f$ is \textbf{surjective} if \\
\begin{center} $\forall \ y\  \in \  Y\  \exists \  x\  \in \  X$ such that $f(x) = y$ \end{center}
This is also called an \textbf{onto} map.
\\
\textbf{Def}: $f$ is \textbf{bijective} if it is bother injective and surjective

\medskip

\begin{enumerate}
	\item{Draw a picture to illustrate a map which is injective but not surjective} 
	\item{Draw a map which is surjective but not injective}
	\item{Prove that a function would have to be bijective in order to be invertible.}
	Proof outline: Consider what happens to the domain and range of a function (including all points within them) when it is inverted. Now consider why invertibility will fail if a function is not injective, and then again if it is not surjective.

\end{enumerate}
\bigskip
\centerline{\textbf{Double angle formula}}
\medskip
The sine double angle formula is
\\
\centerline{$sin(2\theta) = 2sin(\theta)cos(\theta)$}
\\
\\
Draw a unit circle and inscribe a triangle such that two points are on the circle, and one point is at the centre of the circle. An illustration is included below. Use the \textbf{Inscribed Angle Theorem} to prove the double angle formula above as follows.
\begin{enumerate}
	\item{Come up with an expression for the area of the triangle (You might choose to use the area of the parallelogram formed by the triangle)}
	\item{Bisect that triangle into two right triangles. Find an expression for the area for these right triangles.}
	\item{Note that the area of the large triangle must be the same as the area of the sum of the two small triangles.}
\begin{figure}[htbp]
        \centering
        \includegraphics[angle=0,width=1.0\textwidth]{double_angle.jpg}
        \caption{Geometry of proof}
\end{figure}

\end{enumerate}








\end{document}
